
\subsection{Aim 1}
Screening polymeric donor combinatorial libraries for stability by irradiation degradation and green solvent extraction.

Libraries of polymers can be screened for irradiation and water stability by simply irradiating in the pressence of water to degrade poor materials.
Materials that exhibit good stability are likely to be stable.

\subsection{Aim 2}
Tin-free synthesis of small molecule and Y6-like NFAs and polymeric PM6-like donor combinatorial libraries by Knoevenagel condensation in water.

B-ialdehydes and electron poor bi-methylenes can be used to form polymeric and dendritic materials.
These materials can either be built by functional protection and deprotection for small molecule acceptors or via stepwise polymerization for donor materials.
Banzaldehydes can be protected with alcohols, such as ethylene glycol, to form the acetal.
These acetals are base stable and can be deprotected in water with acid.

\subsection{Aim 3}
Implementation of optimized materials into bulk heterojunction OSCs to determine PCE.

\subsection{Potential challenges and solutions}
